% =========================================================
% Down-Sets and Up-Sets
% =========================================================

\section{Down-Sets and Up-Sets}
\label{sec:down-up-sets}

\begin{remark*}[Toolkit: Down-Sets and Up-Sets]
A down-set is a subset of a poset that is closed downward: if an element
belongs to it, everything below that element belongs to it too.  An up-set
is the dual notion --- closed upward.  Down-sets are the order ideals of
lattice theory; up-sets are the order filters.  The collection of all
down-sets of a poset, ordered by inclusion, is itself an important poset.
\end{remark*}

% ---------------------------------------------------------
% def:down-set
% ---------------------------------------------------------

\begin{definitionbox}{Definition (Down-Set)}
\begin{definition}[Down-Set]
\label{def:down-set}
Let $(P, \leq)$ be a partially ordered set.  A subset $Q \subseteq P$ is
a \emph{down-set} (or \emph{decreasing set}, or \emph{order ideal}) if it
is closed downward:
\[
\forall x \in Q,\; \forall y \in P,\quad y \leq x \;\Longrightarrow\; y \in Q.
\]
\end{definition}
\end{definitionbox}

\begin{remark*}[Standard quantified statement]
\[
\operatorname{DownSet}(Q, P, \leq)
\;\coloneqq\;
Q \subseteq P
\;\land\;
\forall x \in Q,\; \forall y \in P,\;
y \leq x \Rightarrow y \in Q.
\]
\end{remark*}

\begin{remark*}[Negated quantified statement]
\[
\neg\operatorname{DownSet}(Q, P, \leq)
\;\coloneqq\;
\exists x \in Q,\; \exists y \in P,\;
y \leq x \land y \notin Q.
\]
\end{remark*}

\begin{remark*}[Interpretation]
A down-set is a downward-closed region of the poset: once an element is
included, the entire lower section below it is forced in.  Down-sets
generalise downward-closed intervals.  The collection of all down-sets of
$P$, denoted $\mathcal{O}(P)$ and ordered by inclusion, is itself a poset
and in the finite case a distributive lattice.  Birkhoff's representation
theorem states that every finite distributive lattice arises this way.
\end{remark*}

\begin{dependencies}
\hyperref[def:lattice-order-partial-order]{Partial order}.
\end{dependencies}

% ---------------------------------------------------------
% def:up-set
% ---------------------------------------------------------

\begin{definitionbox}{Definition (Up-Set)}
\begin{definition}[Up-Set]
\label{def:up-set}
Let $(P, \leq)$ be a partially ordered set.  A subset $Q \subseteq P$ is
an \emph{up-set} (or \emph{increasing set}, or \emph{order filter}) if it
is closed upward:
\[
\forall x \in Q,\; \forall y \in P,\quad x \leq y \;\Longrightarrow\; y \in Q.
\]
\end{definition}
\end{definitionbox}

\begin{remark*}[Standard quantified statement]
\[
\operatorname{UpSet}(Q, P, \leq)
\;\coloneqq\;
Q \subseteq P
\;\land\;
\forall x \in Q,\; \forall y \in P,\;
x \leq y \Rightarrow y \in Q.
\]
\end{remark*}

\begin{remark*}[Negated quantified statement]
\[
\neg\operatorname{UpSet}(Q, P, \leq)
\;\coloneqq\;
\exists x \in Q,\; \exists y \in P,\;
x \leq y \land y \notin Q.
\]
\end{remark*}

\begin{remark*}[Interpretation]
Dual to down-set via the
\hyperref[prop:duality-principle]{Duality Principle}: the up-sets of $P$
are exactly the down-sets of $P^{\mathrm{op}}$.  In topology, filters are
special up-sets closed under finite intersection and satisfying a
directedness condition.  The complement of a down-set is an up-set, and
vice versa.
\end{remark*}

\begin{dependencies}
\hyperref[def:down-set]{Down-set},
\hyperref[prop:duality-principle]{duality principle}.
\end{dependencies}
